\chapter{Introduction} \label{chap:intro}

\section{Clinical Decision Support and the MIMIC-IV Foundation}
\label{sec:cds_mimic_foundation}

Modern intensive care units generate vast quantities of clinical data, from continuous physiological monitoring and laboratory results to unstructured physician notes and discharge summaries. Translating this data deluge into timely, actionable clinical decisions remains one of the central challenges in critical care informatics. Clinical decision support (CDS) systems powered by machine learning offer a promising path forward, enabling automated pattern recognition across high-dimensional feature spaces that exceed the capacity of traditional severity scores and rule-based heuristics.

A key enabler of reproducible CDS research is the \emph{Medical Information Mart for Intensive Care} (MIMIC-IV), a freely accessible, de-identified database covering over 300,000 hospital admissions at the Beth Israel Deaconess Medical Center \citep{Johnson_2023_MIMIC,johnson2023mimicivscidata}. MIMIC-IV provides structured tables for demographics, diagnoses (ICD-9/10-CM codes), procedures, laboratory measurements, medications, and charted observations, alongside free-text clinical notes including discharge summaries. Its breadth and depth make it a natural testbed for developing and benchmarking machine learning pipelines across diverse clinical tasks.

This thesis investigates two complementary clinical decision support problems that share MIMIC-IV as their common data foundation:

\begin{itemize}
    \item \textbf{Problem~1: Automated ICD Code Prediction for Rare Diseases.} Accurate assignment of International Classification of Diseases (ICD) codes from discharge summaries is essential for billing, epidemiological surveillance, and downstream clinical research. The long-tail distribution of ICD codes means that rare conditions are chronically under-represented in training data, leading to poor recall for precisely those diagnoses where coding errors carry the greatest clinical consequence. This problem was investigated through literature review, problem formulation, dataset exploration, and preliminary modeling, but full model training and synthetic data generation were not completed due to high computational infrastructure costs.
    \item \textbf{Problem~2: Early Mortality Risk Stratification for SA-AKI.} Sepsis-associated acute kidney injury (SA-AKI) couples systemic infection with acute renal failure and is associated with sharply elevated in-hospital mortality. Timely identification of high-risk patients within the first 24 hours of ICU admission can guide escalation decisions, nephroprotective strategies, and resource allocation. This problem was fully executed, yielding a leakage-aware predictive pipeline with interpretable classification and survival analysis models.
\end{itemize}

The sections that follow introduce each problem in turn, beginning with the ICD coding challenge and then presenting the SA-AKI mortality prediction task in detail.

\section{Problem 1: Automated ICD Code Prediction for Rare Diseases}
\label{sec:p1:icd_intro}

The \emph{International Classification of Diseases (ICD)} is a globally recognized coding system maintained by the World Health Organization \citep{who2019icd11}, with its most recent iterations being ICD-10 and ICD-11. This system was developed to standardize the labeling of diseases, medical procedures, and underlying health conditions, thereby fostering interoperability across diverse healthcare environments. As a result, patient diagnoses and treatments are recorded more consistently, allowing healthcare providers to communicate more effectively, researchers to examine data with fewer ambiguities, and public health authorities to craft evidence-based policies.

When multiple ICD codes are assigned to a single patient's record, the information becomes richer by including primary diagnoses and related comorbidities. Large-scale databases, such as MIMIC-IV, reflect this practice by covering not only the chief complaint but also secondary or rare conditions \citep{johnson2023mimicivscidata}. A major concern, however, arises if any of these codes are misassigned. Inaccuracies and omissions can ripple through subsequent clinical decision-making, research conclusions, and billing outcomes. Hence, the role of accurate coding becomes paramount in bridging the gap between raw clinical documentation and reliable medical insights.

From a financial standpoint, healthcare institutions depend on correct coding to claim reimbursements, and errors---either over-coding or under-coding---risk financial imbalances. Studies have highlighted the direct impact of clinical documentation on billing accuracy and institutional revenue, stressing the importance of precision in ICD assignments \citep{dong2022automated}. At a population level, aggregated ICD-coded data forms the bedrock of epidemiological tracking, policy development, and larger-scale medical research. Erroneous codes can alter the trajectory of these studies, potentially skewing public health interventions \citep{edin2023automated}.

Nevertheless, the manual process of ICD coding remains both tedious and vulnerable to human error. Coders routinely parse extensive and often unstructured discharge summaries that include technical jargon, acronyms, or shorthand specific to a given clinical context \citep{wrenn2010quantifying}. Rare or lesser-known conditions risk being misrepresented or overlooked in the daily rush of high-volume coding workloads. Automated or semi-automated tools have therefore been suggested as an avenue for alleviating these problems. Such solutions provide preliminary code assignments that human coders can validate, thus streamlining workflow and improving data quality \citep{campbell2020computer}.

Over the years, numerous computational methods have been devised to assist in the assignment of ICD codes. Early rule-based systems depended on dictionaries and handcrafted rules \citep{farkas2008automatic}, but these were limited in scalability and flexibility. Traditional machine learning techniques (e.g., support vector machines and logistic regression) made modest progress but often stumbled over the diverse, nuanced terminology of clinical text \citep{scheurwegs2017data,perotte2014diagnosis}. Momentum grew significantly with the advent of deep learning, starting with convolutional neural networks (CNNs) \citep{kim2014convolutional,mullenbach2018explainable} and recurrent neural networks (e.g., LSTMs \citep{hochreiter1997long}) that learned hierarchical or temporal patterns from unstructured notes. More recently, transformer-based models such as Bidirectional Encoder Representations from Transformers (BERT) have improved performance in many language tasks \citep{devlin2019bert}, though they often struggle with extremely long documents and heavily imbalanced data \citep{gao2021limitations,huang2022plm}.

Several specialized strategies have since been proposed to address these challenges, including extended context transformers, hierarchical embeddings, and curriculum learning over the ICD taxonomy \citep{nickel2017poincare,ren2022hicu,bengio2009curriculum}. These refinements aim to capture relevant details in lengthy clinical narratives and to handle the long-tail distribution in which a few codes dominate while a majority of rare codes receive minimal coverage \citep{rios2018few}. Alongside these methodological expansions, interpretability is receiving growing attention in automated ICD coding: clinicians must verify or scrutinize a model's reasoning to ensure that each assigned code is justified by actual clinical evidence \citep{holzinger2017we,warner2024modernbert}.

In parallel, new lines of research have begun focusing on entity recognition and ontology-based approaches to enhance the robustness of coding systems. Tools like \emph{cTAKES} \citep{savova2010ctakes} and \emph{MedCAT} \citep{kraljevic2021multi} identify medical concepts in text and map them to standardized vocabularies, such as SNOMED~CT \citep{donnelly2006snomedct} or the Unified Medical Language System (UMLS) \citep{bodenreider2004umls}. Although these pipelines do not automatically produce final ICD codes, they offer structured terminological insights that can reduce ambiguity in free-text data \citep{wu2018semehr}. More advanced systems incorporate ICD-specific knowledge directly into the classification process (e.g., hyperbolic embeddings of the ICD hierarchy \citep{cao2020hypercore} or label descriptions drawn from official manuals). This knowledge-aware approach has consistently shown promise in bridging the gap between unstructured narratives and the highly specialized code sets used in clinical practice.

An especially difficult hurdle is how to improve the recognition of rare or ``zebra'' codes. When diagnoses occur infrequently, data-driven models often lack sufficient positive examples for generalization \citep{rios2018few}. To combat this scarcity, researchers have begun to explore \emph{synthetic data generation}. By leveraging large language models (LLMs) to create artificial discharge summaries or augment existing text, the model's exposure to underrepresented conditions can be increased \citep{falis2022horsestozebras,falis2024gpt35}. Ontologies like SNOMED~CT or the UMLS help ground these generated examples so that the synthetic narratives remain medically plausible \citep{kumichev2024medsyn}. For instance, a rare genetic disorder might have very few real cases in a training corpus; generating extra samples keyed to that condition's known symptomatology can yield substantial performance gains, particularly in recall for low-frequency labels.

These generative methods build on a broader trend in natural language processing, in which powerful transformers (e.g., GPT-3.5 or GPT-4) are used to produce text that mimics human writing. Indeed, recent work suggests that LLMs prompted with curated medical knowledge can synthesize relatively realistic patient histories \citep{yang2023fewshoticd}. Although concerns persist about potential hallucinations or a lack of full clinical richness, experiments have shown that carefully designed prompts---possibly incorporating knowledge graphs or code descriptions---can yield synthetic notes beneficial for model training. Complementary approaches use synonyms or sibling diagnoses from an ontology to systematically ``perturb'' or augment existing text \citep{yuan2022codesynonyms}, thus expanding the model's exposure to variant language for the same condition.

However, practical integration of these advanced methods requires consideration of real-world constraints. Data privacy regulations often limit the availability of large, richly annotated clinical corpora. Synthetic notes, if generated responsibly, can alleviate privacy concerns while still reflecting key features of real patient reports. Yet computational cost and resource constraints remain a challenge: training or fine-tuning long-context transformers can be expensive, and using LLM APIs extensively may be prohibitively costly for some healthcare institutions \citep{ji2021bert}. Techniques such as quantization, gradient checkpointing, or hierarchical batching \citep{geiping2023cramming} can reduce resource demands. Another paramount concern is reliability: automated coding systems must provide transparent justifications for the codes they produce if they are to earn the trust of clinicians. Interpretable mechanisms, such as label-specific attention or post-hoc attribution methods (e.g., Integrated Gradients, SHAP), facilitate manual review and verification \citep{teng2020explainable}.

To address these persistent challenges, this thesis proposed a framework that combines ontology-informed data augmentation, cutting-edge neural architectures, and interpretability techniques to enhance automated ICD coding. In particular, four major strategies anchor the work:

\begin{itemize}
    \item \textbf{Synthetic Data Generation for Rare Codes:} Large language models are leveraged in tandem with medical ontologies (e.g., SNOMED~CT, UMLS) to produce realistic training examples for less frequently occurring conditions \citep{kumichev2024medsyn,falis2022horsestozebras}. By enriching the dataset with synthetic discharge summaries, the aim is to close the performance gap between common and rare codes.

    \item \textbf{Knowledge-Aware Multi-Label Architectures:} Neural models capable of handling thousands of code labels are examined, including transformer-based classifiers with extended input contexts. Techniques such as hyperbolic embeddings of the ICD hierarchy \citep{nickel2017poincare,cao2020hypercore} and label attention mechanisms \citep{vu2020ijcai} are integrated to improve performance and interpretability in multi-label classification.

    \item \textbf{Clinically Oriented Interpretability:} Interpretability remains crucial for adoption in clinical workflows. Methods like SHapley Additive exPlanations (SHAP), attention heatmaps, or concept-level attributions \citep{holzinger2017we,teng2020explainable} are incorporated to help coders and physicians quickly validate or question the system's logic, especially for critical or rare diagnoses.

    \item \textbf{Resource-Conscious Training and Deployment:} Training strategies such as quantization and gradient checkpointing are used to mitigate heavy computational requirements \citep{geiping2023cramming}. This ensures that the proposed system remains accessible to healthcare organizations with limited GPU resources.
\end{itemize}

\noindent\textbf{Note on completion status.} The ICD code prediction work (Problem~1) progressed through literature review, formal problem definition, dataset exploration, and preliminary baseline modeling. However, the full synthetic data generation pipeline and large-scale model training were not completed due to the high computational infrastructure costs involved. The methodology and preliminary findings are documented in subsequent chapters, and the remaining steps are outlined as future work in Chapter~\ref{chap:conclusion}.

\bigskip
\noindent Having introduced the ICD coding challenge, the remainder of this chapter focuses on the second problem investigated in this thesis: early mortality prediction for sepsis-associated acute kidney injury (SA-AKI). While Problem~1 addresses the textual domain of clinical documentation, Problem~2 operates on structured clinical variables, and both share MIMIC-IV as their common data foundation.

\section{Problem 2: Clinical Background and Motivation} Sepsis-associated acute kidney injury (SA--AKI) has emerged as one of the most consequential syndromes in critical care because it couples the systemic instability of sepsis with acute loss of renal function and thereby compounds risk across organ systems. For clinical and research purposes SA--AKI is commonly labeled when acute kidney injury develops within seven days of sepsis onset so that kidney dysfunction can be attributed to the septic process rather than to unrelated peri-ICU events \citep{Peerapornratana_2019}. The constituent conditions are defined by international consensus. Sepsis is defined by Sepsis--3 as life-threatening organ dysfunction due to a dysregulated host response to infection and is operationalized by an acute rise in SOFA score of two points or more \citep{Singer_2016}. Acute kidney injury is defined and staged by KDIGO using serum creatinine and urine output thresholds, including an absolute creatinine increase of at least $0.3\,\mathrm{mg/dL}$ within forty-eight hours, a relative increase of at least $1.5\times$ baseline within seven days, or oliguria at or below $0.5\,\mathrm{mL/kg/h}$ for six hours \citep{KDIGO_2012}. Within intensive care units SA--AKI has been reported to account for a large share of AKI cases and to be associated with sharply higher in-hospital mortality, longer length of stay, and increased risk of subsequent chronic kidney disease in survivors \citep{Poston_Koyner_2019,Peerapornratana_2019}. These epidemiological features have kept early risk stratification at the forefront of research because timely identification of clinical deterioration is central to appropriate escalation, nephroprotective strategies, and resource planning.

\section{What the Literature Has Tried and Where It Falls Short} Risk assessment in the ICU has traditionally relied on general severity scores such as SOFA, SAPS~II, and APACHE~II. These scores have offered useful global snapshots of illness severity but they tend to show only fair discrimination when the target is SA--AKI mortality, and they do not represent the nonlinear and interacting physiology that links infection biology, hemodynamics, renal injury, and treatment exposure in this population \citep{Luo_2022,Qu_2024}. In recent years the availability of large electronic health record repositories, notably MIMIC--IV, has encouraged the development of machine learning models that can take advantage of richer feature spaces and flexible function classes \citep{Johnson_2023_MIMIC}. Tree-based ensembles such as extreme gradient boosting frequently outperform baseline scores in internal testing and AUROC values near or above 0.80 are often reported \citep{Chen_2025,Qu_2024,Wang_2024}. Nevertheless, several methodological patterns recur across studies and create uncertainty about the true clinical utility of published models. Temporal anchoring of predictors is sometimes vague and whole-stay aggregation has been used, which allows future information to flow into features that are later presented as early predictors. Proxy variables for time to event, including length of stay, have been included as inputs to mortality models, which introduces a direct leakage channel from outcome process to covariates and inflates apparent discrimination \citep{Chen_2025}. Data-dependent steps such as feature selection and imputation have sometimes been performed across the full dataset rather than being learned on development folds and then applied to held-out data, which again invites optimistic estimates \citep{Kang_Yoon_2023}. Class imbalance has often been handled through synthetic oversampling, including SMOTE and ADASYN, without careful assessment of whether the generated examples concentrate in overlapped or noisy regions or near outliers where decision boundaries are unstable \citep{Chawla_2002,He_2008,Blagus_2013,Saez_2015,Krawczyk_2016}. Calibration and principled threshold selection have been inconsistently reported, which limits bedside interpretability even when AUROC values appear strong. Finally, most model development has relied on data from a single tertiary-care center, and performance losses under transport to other hospitals and populations have been documented, which tempers claims of broad generalizability \citep{Johnson_2023_MIMIC,Rockenschaub_2024}.

\section{Problem Focus and Research Objectives} The central challenge addressed in this thesis is the chasm between the promising internal accuracy reported in machine learning literature and the practical bedside utility of predictive models for SA--AKI mortality. This gap is primarily shaped by two issues: first, a failure to align predictive pipelines with the real-world clinical decision timeline, and second, a lack of transparency and calibration that hinders a model's adoption. The objective of this work is to develop and validate a mortality risk stratification framework for SA--AKI that is methodologically sound, clinically relevant, and transparent. The work is conducted using the MIMIC--IV dataset \citep{Johnson_2023_MIMIC}, with a strict focus on an early clinical horizon to ensure the model's predictions are actionable. The specific aims are: \begin{enumerate} \item To perform a systematic exploratory data analysis to identify key clinical phenotypes and risk factors within the training cohort. \item To develop and validate a well-calibrated LightGBM classification model for in-hospital mortality using only features available within the first 24 hours of ICU admission. \item To ensure model transparency by providing both global and local feature-level explanations using SHapley Additive exPlanations (SHAP) \citep{Lundberg_Lee_2017}. \item To complement the binary prediction with a dynamic risk assessment by developing survival models using Cox Proportional Hazards and DeepSurv to estimate the \emph{time to in-hospital death}. \end{enumerate}

\section{Proposed Approach and Evaluation Plan} The approach adopted in this thesis begins with the disciplined construction of an SA--AKI cohort according to Sepsis--3 and KDIGO criteria, with all predictors strictly limited to the first 24 hours (T0--T24) of an ICU stay \citep{Singer_2016,KDIGO_2012}. This work is distinguished by a leakage-aware preprocessing pipeline where all steps were fitted exclusively on training data. This includes a formal test for the missing data mechanism, followed by the creation and statistical selection of informative missingness indicators and median imputation. A broad benchmark of machine learning models was conducted, leading to the selection of a final LightGBM model based on its superior performance. This model was tuned using the Optuna framework, explicitly calibrated using isotonic regression, and an optimal decision threshold was determined via Youden's J statistic on a validation set. Explanations for the final model's predictions were generated using SHAP to ensure transparency. Finally, to provide a dynamic view of risk over time, two survival models, a traditional Cox Proportional Hazards model and a neural network-based DeepSurv model, were developed and evaluated using the same early-window covariates. The evaluation plan prioritizes methodological transparency and clinical relevance, reporting a suite of metrics including AUROC, F1-score, Harrell's C-index, and calibration plots on a strictly held-out test set, thereby addressing the common pitfalls identified in the literature \citep{Kang_Yoon_2023,Chen_2025}.

\section{Chapter Roadmap} The remainder of this thesis is organized into six chapters that address both research problems in a unified structure.

Chapter~\ref{chap:litreview} reviews the literature in two parts: Part~I surveys the evolution of automated ICD coding, covering rule-based systems, deep learning approaches, transformer architectures, curriculum learning, synthetic data generation for rare codes, and evaluation benchmarks; Part~II reviews SA--AKI epidemiology and definitions, conventional severity scores, machine learning approaches for mortality prediction, and synthesizes methodological lessons from recent studies. A closing synthesis section ties both bodies of work together under the clinical decision support theme.

Chapter~\ref{chap:problem} formally defines both research problems. Problem~1 presents the automated ICD-10-CM multi-label classification task with its mathematical formulation, rare-code challenges, and proposed evaluation framework. Problem~2 states the SA--AKI early mortality prediction task, including the clinical decision scope, observational unit definitions, the three sub-problems (phenotyping, supervised classification, and survival analysis), and the evaluation protocol.

Chapter~\ref{chap:methods} describes the methodology for both problems. The Problem~1 methodology (proposed but not fully executed) covers ontology-based disease grouping, synthetic discharge summary generation, deep learning model training, and MIMIC-IV processing. The Problem~2 methodology details the cohort construction, leakage-aware preprocessing pipeline, model development and calibration process, survival analysis, and adherence to TRIPOD reporting guidelines.

Chapter~\ref{chap:eda} presents the dataset understanding and exploratory data analysis. Part~I covers the MIMIC-IV database from the ICD coding perspective, including SQL query results, visual explorations, and data preparation. Part~II presents the SA--AKI cohort analysis, including source tables, ontology-driven concept identification, cohort construction, clinical labeling, and phenotyping of survivor versus non-survivor groups.

Chapter~\ref{chap:results} reports the results for both problems. Problem~1 presents partial results including the KNN baseline, MinHash similarity detection, and the rare disease co-occurrence pipeline. Problem~2 presents the full results for the classification and survival models on the held-out test set, including discrimination, calibration, SHAP-based interpretations, and comparison with prior work.

Finally, Chapter~\ref{chap:conclusion} concludes with a summary of findings for both problems, discusses clinical implications, and outlines future research directions including the remaining steps for completing the ICD coding pipeline and extensions to the SA--AKI mortality prediction framework.
